\section{Proofs for Finite-Sweep Accuracy}\label{app:finite-sweeps51}
\begin{proof}[Proof of Theorem~\ref{M-thm:finite-sweeps51}]
Fix a pass, date, and state, and suppress their indices. Write $J$ for the
incumbent continuation, $h$ for its approximator, $\pi$ for the incumbent
action, and $v$ for the proposal. By the signed evaluation band, for every
feasible action $a$,
\[
 \beta a_{t+1}\leq Q(J;x,a)-Q(h;x,a)\leq\beta b_{t+1}.
\]
Consequently, the incumbent-relative advantage of the proposal is at most
$\overline d+\beta w_{t+1}$. An accepted proposal therefore does not raise
the incumbent-continuation Bellman cost, while rejection retains exactly that
cost. Backward induction gives $J^{\pi^{k+1}}\leq J^{\pi^k}$ at every
date. This induction uses the new continuation for the implemented new
policy; the comparison itself uses the old continuation.

For an accepted proposal, approximate greediness and the band imply
\[
 Q(J;x,v)\leq (T_th)(x)+\alpha+\beta b_{t+1}
 \leq (T_tJ)(x)+\alpha+\beta w_{t+1}.
\]
For a rejected proposal, $\overline d+\beta w_{t+1}>0$ and
$\overline d\leq d+\zeta$. Thus
\[
 Q(h;x,\pi)<Q(h;x,v)+\zeta+\beta w_{t+1}
 \leq(T_th)(x)+\alpha+\zeta+\beta w_{t+1}.
\]
Transferring this inequality from $h$ to $J$ adds at most one further
$\beta w_{t+1}$. Both cases therefore yield
\begin{equation}\label{eq:approx-greedy-old51}
 T_t^{\pi^{k+1}}J_{t+1}^{\pi^k}
 \leq T_tJ_{t+1}^{\pi^k}+\epsilon_t^k.
\end{equation}
The already proved policy nonincrease, monotonicity of
$T_t^{\pi^{k+1}}$, and \eqref{eq:approx-greedy-old51} give
\[
 J_t^{\pi^{k+1}}\leq T_t^{\pi^{k+1}}J_{t+1}^{\pi^k}
 \leq T_tJ_{t+1}^{\pi^k}+\epsilon_t^k.
\]
The Bellman operator is $\beta_t$-Lipschitz in the supremum norm. Subtracting
$V_t^*=T_tV_{t+1}^*$ and using feasibility proves
\eqref{M-eq:sweep-recursion51}. Repeated substitution proves
\eqref{M-eq:sweep-unroll51}; at $m=T-t$ its terminal term vanishes.
Nonincrease of $E_t^k$ in further passes then gives
\eqref{M-eq:sweep-final51}. The exact case has zero error at every date.
\end{proof}

\begin{proof}[Proof of Proposition~\ref{M-prop:cell-greedy51}]
Fix $x\in C$ and any feasible $a\in A_t(x)$. Choose $m\in M_C$ with
$\|a-m\|\leq\delta_{A,C}$. Then
\begin{align*}
 Q_t(h;x,v_C)&\leq Q_t(h;x_C,v_C)+L_tr_C\\
 &\leq Q_t(h;x_C,m)+\eta_C+L_tr_C\\
 &\leq Q_t(h;x,a)+\eta_C+2L_tr_C+D_t\delta_{A,C}.
\end{align*}
Taking an infimum over $a$ proves the assertion without assuming that the
infimum is attained. Menu feasibility is uniform on $C$ by hypothesis.
\end{proof}

\begin{proof}[Proof of Proposition~\ref{M-prop:sharp-width51}]
Take $T=1$, zero stage cost, $\beta_0=1$, and two deterministic actions
leading to different terminal states $y_0,y_1$. Let $w>0$ and
$0<\delta<w$. Set
\[
 h_1(y_0)=w,\quad h_1(y_1)=\delta,\qquad
 g(y_0)=2w,\quad g(y_1)=\delta.
\]
The exact terminal error lies in $[0,w]$. The incumbent chooses $y_0$ and
the exact $h_1$-greedy proposal chooses $y_1$. Its exact contrast is
$d=\delta-w$, with $\alpha=\zeta=0$. The safe upper comparison is
$d+w=\delta>0$, so the gate retains the incumbent. Its actual optimality
gap is $2w-\delta$. Sending $\delta$ down to zero establishes sharpness.
This example also shows why rejecting an uncertified proposal is not evidence
that the proposal lacks a true economic benefit.
\end{proof}

\paragraph{Computed verification and logical scope.}
The source-bound exact-rational tests use finite transition matrices, signed
residual propagation, and independently evaluated policy and optimal costs.
They verify policy nonincrease, both gate cases, the one-pass recurrence,
constant-shift invariance, finite-horizon exact termination, and the two-action
sharpness construction. They are regression tests for the implementation of
the displayed formulas; the preceding arguments, rather than enumeration,
prove the theorems for the stated model class. When bands are statistical,
all passes and proposed comparisons require a common simultaneous coverage
event or a valid sequential error allocation. Reusing a confidence interval
outside its stated family does not satisfy the hypotheses.
